% !TEX program = pdflatex
\documentclass[10pt]{article}

\usepackage[margin=0.85in]{geometry}
\usepackage{amsmath,amssymb,amsthm}
\usepackage{booktabs}
\usepackage{graphicx}
\usepackage{microtype}
\usepackage{xcolor}
\usepackage[hidelinks]{hyperref}

\newcommand{\KL}{D_{\mathrm{KL}}}
\newcommand{\R}{\mathbb{R}}
\newcommand{\A}{\mathcal{A}}
\newcommand{\E}{\mathbb{E}}
\newcommand{\method}{Online RC-OMD}
\newtheorem{proposition}{Proposition}

\title{Adaptive Credit-Weighted Online Mirror Descent\\
for Sparse Verifiable Reward Reinforcement Learning}
\author{RC-OMD Project Team}
\date{Working draft --- 8 August 2026}

\begin{document}
\maketitle

\begin{abstract}
Sparse terminal rewards make it difficult to decide which actions in a long
trajectory deserve a large policy update.  Broadcasting one group-relative
advantage to every decision is inexpensive, but finite-sample correlations can
move policy mass at decisions that do not affect the reward.  We study a
controlled sequence setting in which the reward-relevant positions are known
for evaluation but hidden from the optimizer.  We introduce Online RC-OMD, a
low-cost reliability-calibrated variant of Online Mirror Descent (OMD).  It
tracks exponentially weighted first and second moments of group-relative
action-score estimates and uses their persistent signal-to-uncertainty ratio to
scale the local KL-mirror step.  The method requires one score computation per
group and $O(H|\A|)$ persistent state.  In a frozen out-of-distribution protocol
covering four threshold-reward tasks and ten seeds per task, Online RC-OMD
remained within $0.0018$ normalized success AUC of a tuned uniform Group OMD
baseline while using only $18.9\%$--$24.8\%$ as much absolute cumulative KL at
known distractor positions.  These results support a narrow Pareto claim about
controlled policy movement; they do not establish causal credit assignment or
performance under function approximation.
\end{abstract}

\section{Introduction}

Reinforcement learning with verifiable rewards (RLVR) supplies an attractive
training signal for tasks whose final outcomes can be checked automatically.
The reward is often binary and observed only after a complete trajectory.  A
group-relative baseline can avoid learning a separate critic, as in Group
Relative Policy Optimization (GRPO)~\cite{shao2024deepseekmath}, but the same
trajectory-level advantage is then commonly propagated through many decisions.
This leaves an optimizer-level question: when the available credit signal is a
noisy proxy, how much policy movement should each decision be allowed to make?

Recent work has attacked token-level credit assignment using entropy-aware
weights and related local signals~\cite{he2026eapo,xie2026acpo}.  Entropy can be
informative, but uncertainty is not identical to reward relevance.  A decision
can be uncertain and irrelevant, or confident and pivotal.  We therefore treat
credit reliability as a quantity to be estimated rather than equating it with
entropy.  Our approach is deliberately smaller than language-model RL: a
factorized tabular policy lets us observe local policy movement exactly and
construct environments with known distractor decisions.

Our contributions are:
\begin{enumerate}
    \item We construct controlled sparse terminal-reward sequence tasks with
    exact counterfactual credit and known reward-irrelevant positions for
    evaluation.  The suite includes settings where policy entropy is aligned
    with, and deliberately misleading about, true positional importance.
    \item We propose \method, which converts persistent agreement in online
    group-relative action scores into per-position OMD step scales.  Unlike a
    per-group bootstrap estimator, it adds no rollout resampling and maintains
    only $O(H|\A|)$ state.
    \item We report development results separately from a protocol frozen
    before evaluation.  The pre-registered comparison passes all four unseen
    scenarios under a declared AUC--distractor-KL rule.
\end{enumerate}

The central empirical claim is a trade-off, not universal dominance: reliability
calibration directs substantially less policy movement toward known distractors
while retaining almost all learning efficiency in the tested controlled tasks.

\section{Related Work}

\paragraph{Mirror-descent policy optimization.}
Mirror descent replaces Euclidean proximity with a Bregman divergence and is
particularly natural on a probability simplex~\cite{beck2003mirror}.  MDPO
connects this geometry to practical trust-region policy optimization by
combining a linearized RL objective with a proximity term~\cite{tomar2020mdpo}.
We use the negative-entropy mirror map, for which the proximity term is KL and
the exact local update is exponentiated gradient.

\paragraph{Group-relative optimization and sparse verifiable rewards.}
DeepSeekMath introduced GRPO as a critic-free relative policy optimization
method for mathematical reasoning~\cite{shao2024deepseekmath}.  Outcome-only
feedback saves the cost of a process reward model, but it leaves temporal credit
coarse.  A recent survey organizes the rapidly growing LLM credit-assignment
literature by granularity and methodology~\cite{zhang2026survey}.

\paragraph{Entropy-aware credit.}
EAPO analyzes reward polarity and token entropy, and shows that entropy bounds
the credit a token can carry under its information-theoretic formulation
\cite{he2026eapo}.  ACPO uses a mode-local surrogate entropy to modulate token
updates asymmetrically~\cite{xie2026acpo}.  Our study does not compete with
these systems at language-model scale.  Instead, it tests a distinct diagnostic
hypothesis: persistent score agreement can regulate local mirror geometry even
when instantaneous entropy is a misleading proxy for reward relevance.

\section{Controlled Problem}

Let a trajectory be $\tau=(a_1,\ldots,a_H)$ with
$a_k\in\A=\{1,\ldots,A\}$.  The factorized policy is
$\pi(\tau)=\prod_{k=1}^{H}\pi_k(a_k)$.  A hidden set
$\mathcal{K}\subseteq\{1,\ldots,H\}$ contains the pivotal positions.  Each
$k\in\mathcal{K}$ has a target action $a_k^\star$.  The terminal reward is
\begin{equation}
R(\tau)=\mathbf{1}\!\left\{
\sum_{k\in\mathcal{K}}\mathbf{1}\{a_k=a_k^\star\}\ge m
\right\}.
\label{eq:reward}
\end{equation}
Positions outside $\mathcal{K}$ are distractors: changing them never changes
the reward.  The special case $m=|\mathcal{K}|$ is an all-match or ``needle''
task.  Threshold tasks use $m<|\mathcal{K}|$.

For each update, $K$ trajectories are sampled and assigned group-relative
advantages
\begin{equation}
\widehat A_i=R(\tau_i)-\frac{1}{K}\sum_{j=1}^{K}R(\tau_j).
\label{eq:group_advantage}
\end{equation}
We estimate the gain for action $a$ at position $k$ with inverse propensity
weighting,
\begin{equation}
\widehat g_{t,k,a}=
\frac{1}{K}\sum_{i=1}^{K}
\frac{\mathbf{1}\{a_{i,k}=a\}\widehat A_i}
{\max(\pi_{t,k}(a),p_{\min})},
\label{eq:score}
\end{equation}
with a fixed importance cap in the implementation.  Uniform Group OMD applies
this noisy action score at every position.  Although the expected distractor
score is zero in our controlled environment, its finite-group realization need
not be, so repeated updates may accumulate irrelevant KL drift.

For diagnosis only, the known transition and reward structure gives exact
on-trajectory counterfactual credit
\begin{equation}
c_k(\tau,\pi)=Q^\pi(s_k,a_k)-V^\pi(s_k).
\end{equation}
The optimizer never receives the pivotal set or oracle credit unless explicitly
labelled as an oracle baseline.

\section{Reliability-Calibrated OMD}

\subsection{Uniform KL-mirror update}

At position $k$, negative-entropy OMD solves
\begin{equation}
\pi_{t+1,k}=\arg\max_{p\in\Delta_A}
\left\{\eta\langle \widehat g_{t,k},p\rangle
-\KL(p\,\|\,\pi_{t,k})\right\},
\end{equation}
with solution
\begin{equation}
\pi_{t+1,k}(a)=
\frac{\pi_{t,k}(a)\exp(\eta\widehat g_{t,k,a})}
{\sum_b\pi_{t,k}(b)\exp(\eta\widehat g_{t,k,b})}.
\label{eq:uniform_omd}
\end{equation}

\subsection{Online reliability estimate}

We center each score vector across actions and maintain exponentially weighted
first and second moments with decay $\beta$.  Let $\widehat\mu_{t,k}$ be the
running mean vector, $\widehat s_{t,k}$ its estimated standard-error vector,
and $n_{\mathrm{eff},t}$ the effective sample size of the exponential weights.
We define
\begin{align}
q_{t,k}
&=\min\left\{1,
\frac{n_{\mathrm{eff},t}}{n_{\mathrm{warm}}}
\right\}
\left[
1-z\frac{\|\widehat s_{t,k}\|_2}
{\max(\|\widehat\mu_{t,k}\|_2,\epsilon)}
\right]_+, \\
\alpha_{t,k}
&=\alpha_{\min}+(1-\alpha_{\min})q_{t,k}.
\label{eq:reliability}
\end{align}
The first factor prevents confidence before enough effective observations; the
second rewards a persistent signal relative to its uncertainty.  The local OMD
update replaces $\eta$ in Eq.~\eqref{eq:uniform_omd} by
$\eta\alpha_{t,k}$.  The current score determines direction, whereas the
historical reliability estimate controls how far the policy may move.

\begin{proposition}[Exact local movement control]
For any actions $a,b$ with positive probability, the reliability-calibrated
update satisfies
\begin{equation}
\log\frac{\pi_{t+1,k}(a)}{\pi_{t+1,k}(b)}
-\log\frac{\pi_{t,k}(a)}{\pi_{t,k}(b)}
=\eta\alpha_{t,k}
\bigl(\widehat g_{t,k,a}-\widehat g_{t,k,b}\bigr).
\end{equation}
Thus reliability scales every local log-odds displacement linearly, without
changing the ranking induced by the current action-score vector.
\end{proposition}

\begin{proof}
Take the log ratio of the two normalized exponentiated-gradient probabilities;
the shared normalizer cancels.
\end{proof}

The estimator uses one action-score computation per group.  Its first and second
moments require $O(HA)$ memory and $O(HA)$ arithmetic per update, matching the
asymptotic score-update cost of the uniform baseline.

\section{Experimental Design}

We separate exploratory development from frozen evaluation.  Development tasks
were used to construct entropy counterexamples, compare bootstrap and online
reliability, and choose the online configuration.  The subsequent OOD protocol,
including tasks, seeds, metrics, parameter values, and decision rule, was
committed before execution under identifier \texttt{ood-v1-2026-08-08}.

The frozen comparison used ten seeds per task.  Uniform Group OMD used base step
$0.75$.  Online RC-OMD used base step $1.25$, decay $0.9$, $z=1$, warm-up $8$,
and reliability floor $0.1$.  These unequal base steps were selected on
development tasks to compare a Pareto-matched operating point; a separate
same-step diagnostic was also frozen.  Normalized success AUC is the mean exact
success probability across evaluation checkpoints.  Distractor KL is the sum
over training iterations and known distractor positions of the absolute local
KL movement.

A scenario passed when Online RC-OMD was no more than $0.01$ below Uniform in
success AUC and used at most $50\%$ of Uniform's distractor KL.  The method-level
Go decision required at least three of four scenario passes.  Runtime feasibility
required an Online-to-Uniform ratio below $1.5$ in at least three scenarios.

\section{Results}

\subsection{Pre-registered OOD evaluation}

\begin{table}[t]
\centering
\small
\caption{Frozen OOD comparison.  AUC difference and KL ratio are Online minus
Uniform and Online divided by Uniform, respectively.  Values are means across
ten seeds; the seed is the replication unit.}
\label{tab:ood}
\begin{tabular}{lrrrrr}
\toprule
Scenario & Uniform AUC & Online AUC & $\Delta$AUC & KL ratio & Runtime ratio\\
\midrule
Dense 2-of-6        & 0.992182 & 0.990865 & -0.001317 & 0.248 & 1.107\\
Needle 5-of-5       & 0.980026 & 0.980782 &  0.000755 & 0.239 & 1.081\\
Threshold 3-of-5    & 0.992195 & 0.990567 & -0.001628 & 0.189 & 1.069\\
Threshold 4-of-6    & 0.989446 & 0.987688 & -0.001758 & 0.199 & 1.079\\
\bottomrule
\end{tabular}
\end{table}

All four scenarios passed the pre-declared rule (Table~\ref{tab:ood}), yielding
a Go decision.  Online RC-OMD retained AUC to within $0.0018$ while using between
$18.9\%$ and $24.8\%$ as much distractor KL.  Runtime overhead was
$6.9\%$--$10.7\%$, so all four scenarios also passed the systems rule.  The
scenarios changed reward structure, horizon, group size, initial success
probability, and included both two- and three-action policies.

\IfFileExists{../results/ood_preregistered/auc_distractor_kl_pareto.png}{
\begin{figure}[t]
\centering
\includegraphics[width=0.82\linewidth]{../results/ood_preregistered/auc_distractor_kl_pareto.png}
\caption{Pre-registered OOD success--distractor-drift comparison.  The generated
artifact is reproduced from the frozen configuration; the numerical decision
is based on the machine-readable summary rather than the figure.}
\label{fig:pareto}
\end{figure}
}{}

At a common base step of $1.0$, Online RC-OMD had $0.0044$--$0.0069$ lower AUC,
but used only $10.7\%$--$13.4\%$ as much distractor KL.  This diagnostic confirms
that the main result is an update-allocation trade-off, not evidence that
reliability scaling always raises reward at an identical nominal step size.

\subsection{Development diagnostics}

In entropy-aligned development tasks, entropy weighting improved early learning;
when distractor positions were initialized with high entropy and pivotal
positions with lower entropy, the same heuristic reduced sample efficiency.
This controlled reversal motivated reliability estimation.  Per-group bootstrap
reliability reduced distractor drift but cost roughly $7$--$9\times$ the uniform
runtime.  Replacing bootstrap resampling with online moments reduced runtime to
within roughly $7\%$--$13\%$ of Uniform on the development tasks.  These
measurements are exploratory and are not pooled with the frozen OOD result.

\section{Discussion and Limitations}

The results are consistent with a simple mechanism: reward-relevant score
directions recur, whereas finite-sample distractor directions fluctuate, so
running agreement can allocate a larger KL budget to the former.  This is an
associational mechanism, not a causal-credit theorem.  A persistent spurious
correlation can also look reliable, and estimator lag can suppress newly useful
directions after a policy or task change.

The current scope has five important limits.  First, the policy is tabular and
factorized; shared neural parameters may couple local updates.  Second, known
distractor positions are used only to evaluate drift, an oracle unavailable in
real RLVR.  Third, the principal methods use base steps chosen on development
tasks.  Fourth, ten seeds quantify simulation variation but not universality.
Fifth, no language model, process reward model, or long natural-language
trajectory has been tested.  These limitations make function approximation the
next necessary experiment, while a language-model study remains optional.

\section{Conclusion}

We studied whether the reliability of a noisy step-level credit proxy should
control local mirror-descent geometry under sparse terminal rewards.  A
low-cost online moment estimator produced a reproducible Pareto improvement in
four frozen controlled OOD scenarios: nearly unchanged success learning with
substantially less policy movement at known distractor decisions.  The next
stage must determine whether this behavior survives shared-parameter function
approximation, where local policy constraints are only approximately realizable.

\begin{thebibliography}{9}

\bibitem{beck2003mirror}
A. Beck and M. Teboulle.
Mirror descent and nonlinear projected subgradient methods for convex
optimization.
\emph{Operations Research Letters}, 31(3):167--175, 2003.

\bibitem{he2026eapo}
Y. He, H. Wu, S. Liu, H. Ge, H. Zhou, K. Wu, Z. Zheng, Q. Lin, Z. Zhong,
and Y. Zhang.
Rethinking token-level credit assignment in RLVR: A polarity-entropy analysis.
\emph{arXiv preprint arXiv:2604.11056}, 2026.

\bibitem{shao2024deepseekmath}
Z. Shao, P. Wang, Q. Zhu, R. Xu, J. Song, X. Bi, H. Zhang, M. Zhang,
Y. K. Li, Y. Wu, and D. Guo.
DeepSeekMath: Pushing the limits of mathematical reasoning in open language
models.
\emph{arXiv preprint arXiv:2402.03300}, 2024.

\bibitem{tomar2020mdpo}
M. Tomar, L. Shani, Y. Efroni, and M. Ghavamzadeh.
Mirror Descent Policy Optimization.
\emph{arXiv preprint arXiv:2005.09814}, 2020.

\bibitem{xie2026acpo}
Z. Xie, Y. You, Y. Li, E. Gong, Z. Chen, Q. Chen, Y. Cheng, P. Jiang,
and Y. Mu.
ACPO: Adaptive credit policy optimization via fine-grained surrogate entropy.
\emph{arXiv preprint arXiv:2607.03126}, 2026.

\bibitem{zhang2026survey}
C. Zhang.
From reasoning to agentic: Credit assignment in reinforcement learning for
large language models.
\emph{arXiv preprint arXiv:2604.09459}, 2026.

\end{thebibliography}

\end{document}
